\chapter{Second-Order Linear ODEs}

Second-order linear ODEs describe springs, masses, pendulums (for small angles), and RLC circuits. They are the most important ODEs in engineering.

\section{Homogeneous equations with constant coefficients}

\begin{defbox}{Standard form}
A second-order linear ODE with constant coefficients is
\[ ay'' + by' + cy = g(x). \]
If $g(x)=0$ the equation is \textbf{homogeneous}. Otherwise it is \textbf{non-homogeneous}. For an IVP we give two conditions: $y(x_0)=y_0$ and $y'(x_0)=y_1$.
\end{defbox}

\begin{methodbox}{Characteristic equation}
To solve $ay'' + by' + cy = 0$, try $y = e^{rx}$. This gives the \textbf{characteristic equation}
\[ ar^2 + br + c = 0 . \]
\begin{center}
\begin{tabular}{@{}lll@{}}
\toprule
Roots & General solution \\
\midrule
Two distinct real roots $r_1, r_2$ & $y = c_1e^{r_1x} + c_2e^{r_2x}$ \\[3pt]
One repeated real root $r$ & $y = (c_1 + c_2x)\,e^{rx}$ \\[3pt]
Complex roots $\alpha \pm i\beta$ & $y = e^{\alpha x}\big(c_1\cos\beta x + c_2\sin\beta x\big)$ \\
\bottomrule
\end{tabular}
\end{center}
\end{methodbox}

\begin{defbox}{Wronskian}
Two solutions $y_1, y_2$ are \textbf{linearly independent} (so they form a basis) if their Wronskian is not zero:
\[ W(y_1,y_2) = y_1y_2' - y_2y_1' \neq 0 . \]
For a linear IVP with continuous coefficients, a unique solution always exists. This is stronger than the first-order case.
\end{defbox}

\begin{derivbox}{The characteristic equation and its three cases}
\textbf{Superposition.} The equation $ay''+by'+cy=0$ is linear and homogeneous. If $y_1$ and $y_2$ are solutions, then $(c_1y_1 + c_2y_2)'' = c_1y_1'' + c_2y_2''$, and so on, so $c_1y_1 + c_2y_2$ is also a solution.

\textbf{The exponential guess.} Try $y = e^{rx}$. Then $y' = re^{rx}$ and $y'' = r^2e^{rx}$, so
\[ (ar^2 + br + c)\,e^{rx} = 0 \ \Longrightarrow\ ar^2 + br + c = 0, \]
because $e^{rx}\neq0$.

\textbf{Case 1: two distinct real roots.} We get two solutions $e^{r_1x}$ and $e^{r_2x}$. Their Wronskian is $W = (r_2 - r_1)e^{(r_1+r_2)x}\neq0$, so they are independent, and superposition gives $y = c_1e^{r_1x} + c_2e^{r_2x}$.

\textbf{Case 2: repeated root.} Here $r = -b/(2a)$ and $e^{rx}$ is only one solution. Look for a second one in the form $y = u(x)e^{rx}$ (``reduction of order''). Then $y' = (u' + ru)e^{rx}$ and $y'' = (u'' + 2ru' + r^2u)e^{rx}$. Substituting,
\[ a u'' + (2ar + b)\,u' + (ar^2 + br + c)\,u = 0 . \]
Since $r$ is a root, $ar^2+br+c = 0$. Since $r = -b/2a$, also $2ar + b = 0$. So $au''=0$, $u = c_1 + c_2x$, and $y = (c_1 + c_2x)e^{rx}$.

\textbf{Case 3: complex roots $\alpha\pm i\beta$.} Euler's formula $e^{i\theta} = \cos\theta + i\sin\theta$ gives the complex solution $e^{(\alpha+i\beta)x} = e^{\alpha x}\cos\beta x + i\,e^{\alpha x}\sin\beta x$. Write it as $y = u + iv$ with $u,v$ real. The coefficients $a,b,c$ are real, so $a y''+by'+cy = (au''+bu'+cu) + i(av''+bv'+cv) = 0$ means the real part and the imaginary part are both zero. So $u = e^{\alpha x}\cos\beta x$ and $v = e^{\alpha x}\sin\beta x$ are two real solutions. Their Wronskian is $\beta e^{2\alpha x}\neq0$.
\end{derivbox}

\begin{examplebox}{Distinct real roots (IVP)}
Solve $y'' + 4y' + 3y = 0$, \; $y(0)=2$, \; $y'(0)=0$.

Characteristic equation: $r^2 + 4r + 3 = 0$, so $r = -1, -3$. General solution: $y = c_1e^{-x} + c_2e^{-3x}$.

Conditions: $c_1 + c_2 = 2$ and $-c_1 - 3c_2 = 0$. So $c_1 = -3c_2$, then $-2c_2 = 2$, giving $c_2=-1$ and $c_1 = 3$:
\[ y = 3e^{-x} - e^{-3x}. \]
\end{examplebox}

\begin{examplebox}{Complex roots (IVP)}
Solve $y'' + 2y' + 5y = 0$, \; $y(0)=1$, \; $y'(0)=0$.

$r^2 + 2r + 5 = 0$ gives $r = -1 \pm 2i$. So $y = e^{-x}(A\cos 2x + B\sin 2x)$.

$y(0) = A = 1$. The derivative is $y' = -e^{-x}(A\cos 2x + B\sin 2x) + e^{-x}(-2A\sin 2x + 2B\cos 2x)$, so $y'(0) = -A + 2B = 0$ and $B = \tfrac12$:
\[ y = e^{-x}\Big(\cos 2x + \tfrac12\sin 2x\Big). \]
This is a decaying oscillation.
\end{examplebox}

\begin{examplebox}{Repeated root (IVP)}
Solve $y'' - 6y' + 9y = 0$, \; $y(0)=1$, \; $y'(0)=2$.

$r^2 - 6r + 9 = (r-3)^2$, so $r=3$ twice and $y = (c_1 + c_2x)e^{3x}$. Then $c_1 = 1$. The derivative is $y' = c_2e^{3x} + 3(c_1 + c_2x)e^{3x}$, so $y'(0) = c_2 + 3 = 2$ and $c_2 = -1$:
\[ y = (1 - x)\,e^{3x}. \]
\end{examplebox}

\section{Non-homogeneous equations}

\begin{defbox}{Structure of the solution}
The general solution of $ay'' + by' + cy = g(x)$ is
\[ y = y_h + y_p, \]
where $y_h$ is the general solution of the homogeneous equation and $y_p$ is \textbf{any one} particular solution of the full equation. \textbf{Apply the initial conditions only after you add $y_h$ and $y_p$.}
\end{defbox}

\begin{derivbox}{Why $y = y_h + y_p$}
Let $y_p$ be one solution of $ay''+by'+cy = g$, and let $y$ be any other solution. The difference $w = y - y_p$ satisfies
\[ aw'' + bw' + cw = g - g = 0, \]
so $w$ is a solution of the homogeneous equation. Hence $y = y_p + w$ with $w = y_h$. Conversely, $y_p + y_h$ is always a solution. So this formula contains all solutions.
\end{derivbox}

\subsection{Undetermined coefficients}

\begin{methodbox}{Undetermined coefficients}
Use it when $g(x)$ is a polynomial, an exponential, a sine or cosine, or a sum or product of these.
\begin{center}
\begin{tabular}{@{}ll@{}}
\toprule
$g(x)$ & Guess for $y_p$ \\
\midrule
$ke^{mx}$ & $Ae^{mx}$ \\
polynomial of degree $n$ & $A_nx^n + \dots + A_1x + A_0$ \\
$k\cos\omega x$ or $k\sin\omega x$ & $A\cos\omega x + B\sin\omega x$ \\
\bottomrule
\end{tabular}
\end{center}
\textbf{Modification rule.} If a term of your guess is already a solution of the homogeneous equation, multiply the whole guess by $x$ (or by $x^2$ if it is a double root).
\end{methodbox}

\begin{examplebox}{Polynomial forcing (IVP)}
Solve $y'' - 3y' + 2y = 4x$, \; $y(0)=1$, \; $y'(0)=0$.

\emph{Homogeneous part:} $r^2 - 3r + 2 = 0$, so $r=1,2$ and $y_h = c_1e^{x} + c_2e^{2x}$.

\emph{Particular part:} try $y_p = Ax + B$. Then $y_p' = A$, $y_p''=0$. Substitute: $-3A + 2Ax + 2B = 4x$. So $A = 2$ and $-6 + 2B = 0$, giving $B=3$. Thus $y_p = 2x + 3$.

\emph{General solution:} $y = c_1e^{x} + c_2e^{2x} + 2x + 3$.

\emph{Conditions:} $y(0) = c_1 + c_2 + 3 = 1$, so $c_1 + c_2 = -2$. $y'(0) = c_1 + 2c_2 + 2 = 0$, so $c_1 + 2c_2 = -2$. Subtract: $c_2 = 0$, $c_1 = -2$.
\[ y = -2e^{x} + 2x + 3. \]
\end{examplebox}

\begin{examplebox}{Modification rule}
Solve $y'' - 2y' + y = e^{x}$.

$y_h = (c_1 + c_2x)e^{x}$. The forcing $e^{x}$ and even $xe^{x}$ are already in $y_h$, so multiply the guess by $x^2$: $y_p = Ax^2e^{x}$.

Compute $y_p' = A(x^2 + 2x)e^{x}$ and $y_p'' = A(x^2 + 4x + 2)e^{x}$. Then
$y_p'' - 2y_p' + y_p = A\,e^{x}\big[(x^2+4x+2) - 2(x^2+2x) + x^2\big] = 2Ae^{x}$.
So $2A = 1$, $A = \tfrac12$, and
\[ y = (c_1 + c_2x)e^{x} + \tfrac12x^2e^{x}. \]
\end{examplebox}

\subsection{Variation of parameters}

Use this when $g(x)$ is not a simple type (for example $\tan x$ or $1/x$).

\begin{methodbox}{Variation of parameters}
For $y'' + p(x)y' + q(x)y = g(x)$ (the coefficient of $y''$ must be 1), let $y_1, y_2$ be two independent solutions of the homogeneous equation and $W = y_1y_2' - y_2y_1'$. Then
\[ y_p = -y_1\int \frac{y_2\,g}{W}\,\dd x + y_2\int \frac{y_1\,g}{W}\,\dd x . \]
\end{methodbox}

\begin{derivbox}{The variation-of-parameters formula}
Start with $y'' + py' + qy = g$ and two independent homogeneous solutions $y_1, y_2$. Look for $y_p = u_1(x)y_1 + u_2(x)y_2$, where the ``constants'' are now functions. We have freedom to choose one condition on $u_1,u_2$. Choose
\[ u_1'y_1 + u_2'y_2 = 0 . \]
Then $y_p' = u_1y_1' + u_2y_2'$ (the extra terms cancel), and
\[ y_p'' = u_1y_1'' + u_2y_2'' + u_1'y_1' + u_2'y_2'. \]
Substitute into the equation and group the terms by $u_1$ and $u_2$:
\[ u_1\big(y_1'' + py_1' + qy_1\big) + u_2\big(y_2'' + py_2' + qy_2\big) + u_1'y_1' + u_2'y_2' = g. \]
The two brackets are zero. So we have two equations for $u_1', u_2'$:
\[ \begin{aligned} u_1'y_1 + u_2'y_2 &= 0,\\ u_1'y_1' + u_2'y_2' &= g. \end{aligned} \]
Solve by Cramer's rule. The determinant is the Wronskian $W = y_1y_2' - y_2y_1'$:
\[ u_1' = -\frac{y_2\,g}{W}, \qquad u_2' = \frac{y_1\,g}{W}. \]
Integrate to get $u_1$, $u_2$, which gives the formula in the method box.
\end{derivbox}

\begin{examplebox}{Variation of parameters}
Solve $y'' + y = \tan x$.

$y_1 = \cos x$, $y_2 = \sin x$, $W = \cos^2x + \sin^2x = 1$. Then
\[ y_p = -\cos x\int \sin x\tan x\,\dd x + \sin x\int \cos x\tan x\,\dd x . \]
Now $\int \sin x\tan x\,\dd x = \int\frac{\sin^2x}{\cos x}\,\dd x = \int(\sec x - \cos x)\,\dd x = \ln|\sec x + \tan x| - \sin x$, and $\int \cos x\tan x\,\dd x = \int\sin x\,\dd x = -\cos x$. So
\[ y_p = -\cos x\ln|\sec x + \tan x| + \sin x\cos x - \sin x\cos x = -\cos x\,\ln|\sec x + \tan x|. \]
The general solution is $y = c_1\cos x + c_2\sin x - \cos x\,\ln|\sec x + \tan x|$.
\end{examplebox}

\section{Euler--Cauchy equations}

\begin{methodbox}{Euler--Cauchy equation}
The equation $x^2y'' + axy' + by = 0$ (for $x>0$) has variable coefficients but is easy. Try $y = x^m$. This gives
\[ m(m-1) + am + b = 0. \]
\begin{center}
\begin{tabular}{@{}ll@{}}
\toprule
Roots & Solution \\
\midrule
distinct real $m_1, m_2$ & $c_1x^{m_1} + c_2x^{m_2}$ \\
repeated $m$ & $(c_1 + c_2\ln x)\,x^{m}$ \\
complex $\alpha\pm i\beta$ & $x^{\alpha}\big(c_1\cos(\beta\ln x) + c_2\sin(\beta\ln x)\big)$ \\
\bottomrule
\end{tabular}
\end{center}
\end{methodbox}

\begin{derivbox}{Euler--Cauchy by a change of variable}
Put $x = e^{t}$ (so $t = \ln x$) and $y(x) = Y(t)$. By the chain rule, $\dfrac{\dd t}{\dd x} = \dfrac1x$, so
\[ y' = \frac{Y'}{x}, \qquad y'' = \frac{\dd}{\dd x}\Big(\frac{Y'}{x}\Big) = \frac{Y''}{x^2} - \frac{Y'}{x^2}. \]
Hence $xy' = Y'$ and $x^2y'' = Y'' - Y'$. The equation $x^2y'' + axy' + by = 0$ becomes
\[ Y'' + (a-1)\,Y' + bY = 0, \]
which has constant coefficients. Its characteristic equation is $m^2 + (a-1)m + b = m(m-1) + am + b = 0$, the same as before. The solutions $e^{mt}$ are $x^m$. A repeated root gives $(c_1 + c_2t)e^{mt} = (c_1 + c_2\ln x)x^m$, and complex roots give $x^\alpha\cos(\beta\ln x)$ and $x^\alpha\sin(\beta\ln x)$.
\end{derivbox}

\begin{examplebox}{Euler--Cauchy}
Solve $x^2y'' - 2xy' + 2y = 0$. Here $m(m-1) - 2m + 2 = m^2 - 3m + 2 = 0$, so $m = 1, 2$ and
\[ y = c_1x + c_2x^2. \]
\end{examplebox}

\section{Modelling: springs and circuits}

\textbf{Mass--spring--damper.} A mass $m$ on a spring of stiffness $k$, with damping coefficient $c$ and external force $F(t)$:
\[ my'' + cy' + ky = F(t). \]
The same equation describes a series RLC circuit: $Lq'' + Rq' + q/C = E(t)$, where $q$ is the charge.

\begin{derivbox}{The spring--mass--damper equation}
Let $y(t)$ be the displacement from equilibrium. Three forces act on the mass $m$:
\begin{itemize}
  \item the spring force (Hooke's law): $-ky$, opposite to the displacement;
  \item the damping force (viscous): $-cy'$, opposite to the velocity;
  \item the external force $F(t)$.
\end{itemize}
Newton's second law $m\,y'' = \text{sum of forces}$ gives $my'' = -ky - cy' + F(t)$, that is
\[ my'' + cy' + ky = F(t). \]
For $F=0$ try $y=e^{rt}$: $mr^2 + cr + k=0$, so
\[ r = \frac{-c\pm\sqrt{c^2 - 4mk}}{2m}. \]
The sign of $c^2-4mk$ decides the case: positive gives two real negative roots (overdamped), zero gives one repeated root (critical), negative gives complex roots with real part $-c/2m<0$ (underdamped). In all three cases $\operatorname{Re} r<0$, so the motion dies out when $c>0$.
\end{derivbox}

\begin{defbox}{Free motion ($F=0$): three cases}
The characteristic equation is $mr^2 + cr + k = 0$.
\begin{itemize}
  \item $c^2 > 4mk$: \textbf{overdamped}. Two negative real roots. The motion returns slowly with no oscillation.
  \item $c^2 = 4mk$: \textbf{critically damped}. A repeated root. The fastest return without oscillation.
  \item $c^2 < 4mk$: \textbf{underdamped}. Complex roots. The motion oscillates with decaying amplitude.
\end{itemize}
Without damping, $c=0$, the natural frequency is $\omega_0 = \sqrt{k/m}$.
\end{defbox}

\begin{examplebox}{Underdamped spring}
A mass $m=1$ kg, damping $c=2$ N\,s/m, and stiffness $k=5$ N/m is pulled $0.1$ m and released from rest. Then $y'' + 2y' + 5y = 0$, $y(0) = 0.1$, $y'(0)=0$.

The roots are $-1\pm 2i$ (underdamped), so $y = e^{-t}(A\cos 2t + B\sin 2t)$. $A = 0.1$ and $-A + 2B = 0$ gives $B = 0.05$:
\[ y = e^{-t}\big(0.1\cos 2t + 0.05\sin 2t\big)\ \text{m}. \]
\end{examplebox}

\begin{derivbox}{Forced undamped oscillation}
Solve $y'' + \omega_0^2y = F_0\cos\omega t$. The homogeneous solution is $c_1\cos\omega_0t + c_2\sin\omega_0t$.

\emph{Case $\omega\neq\omega_0$.} Try $y_p = A\cos\omega t$. Then $y_p'' = -\omega^2A\cos\omega t$, so $A(\omega_0^2 - \omega^2)\cos\omega t = F_0\cos\omega t$ and $A = \dfrac{F_0}{\omega_0^2 - \omega^2}$. This becomes large when $\omega\to\omega_0$.

\emph{Case $\omega = \omega_0$.} The guess $A\cos\omega_0t$ is already a homogeneous solution, so by the modification rule we use $y_p = t\,(A\cos\omega_0t + B\sin\omega_0t)$. Compute $y_p'' = 2\omega_0(-A\sin\omega_0t + B\cos\omega_0t) - \omega_0^2y_p$, so $y_p'' + \omega_0^2y_p = 2\omega_0\big(B\cos\omega_0t - A\sin\omega_0t\big) = F_0\cos\omega_0t$. Hence $A=0$ and $B = \dfrac{F_0}{2\omega_0}$, which gives $y_p = \dfrac{F_0}{2\omega_0}\,t\sin\omega_0t$. The factor $t$ makes the amplitude grow linearly.
\end{derivbox}

\begin{defbox}{Resonance}
For the undamped system $y'' + \omega_0^2y = F_0\cos\omega t$ with $\omega \neq \omega_0$, the particular solution is $y_p = \dfrac{F_0}{\omega_0^2 - \omega^2}\cos\omega t$. The amplitude becomes large when $\omega$ is close to $\omega_0$. When $\omega = \omega_0$ the modification rule gives
\[ y_p = \frac{F_0}{2\omega_0}\,t\sin\omega_0 t, \]
and the amplitude grows without bound. This is \textbf{resonance}. Real systems have damping, so the amplitude stays finite, but it can still be dangerous.
\end{defbox}

\begin{practicebox}
\begin{enumerate}
  \item Solve $y'' - 5y' + 6y = 0$, \; $y(0)=1$, \; $y'(0)=0$.
  \item Solve $y'' + 9y = 0$, \; $y(0)=2$, \; $y'(0)=3$.
  \item Solve $y'' + 2y' + y = 0$, \; $y(0)=0$, \; $y'(0)=1$.
  \item Find the general solution of $y'' + y' - 2y = 6e^{x}$.
  \item Find the general solution of $y'' + 4y = 8x$.
  \item Solve $x^2y'' + 3xy' + y = 0$.
  \item A $2$ kg mass on a spring with $k=8$ N/m has no damping. It starts at $y(0)=0.2$ m with zero velocity. Find $y(t)$.
\end{enumerate}
\end{practicebox}

\begin{answerbox}
(1) $y = 3e^{2x} - 2e^{3x}$. \quad (2) $y = 2\cos 3x + \sin 3x$. \quad (3) $y = xe^{-x}$. \quad (4) $y = c_1e^{x} + c_2e^{-2x} + 2xe^{x}$ (the modification rule is needed because $e^{x}$ is in $y_h$). \quad (5) $y = c_1\cos 2x + c_2\sin 2x + 2x$. \quad (6) $m^2 + 2m + 1 = 0$, so $m=-1$ twice and $y = (c_1 + c_2\ln x)/x$. \quad (7) $\omega_0 = \sqrt{8/2} = 2$ rad/s, so $y = 0.2\cos 2t$ m.
\end{answerbox}
